\section{Central induction and two twist moments}\label{sec:finite-induction}
Let $\mathcal C$ be pseudo-unitary, with its canonical spherical structure,
and let $F:\mathcal Z(\mathcal C)\to\mathcal C$ have right adjoint $I$.
For simples $X_i$ and $A$, write
$M_{iA}=[F(A):X_i]=[I(X_i):A]$. The four simple summands of $I(1)$
occur once, have twist one and dimensions $D/f_\chi$
\cite[Theorems 2.5 and 2.13]{OstrikRank3}.
Let $L$ be their four restriction columns, including the central unit
column $(1,0,0,0)^{\mathsf T}$. Then
\begin{equation}\label{eq:induction-gram}
 MM^{\mathsf T}=B,\qquad
 L\mathbf1=B_{\bullet0},\qquad
 B-LL^{\mathsf T}\text{ is a nonnegative integral Gram matrix}.
\end{equation}
Indeed $FI(X)=\bigoplus_j X_j\otimes X\otimes X_j^*$, and adjunction
gives the first identity. All remaining columns have first entry zero,
which gives the other assertions. In particular the remainder is positive
semidefinite and entrywise nonnegative. A vector in its kernel is
orthogonal to every remaining restriction column.

\subsection{Higher vanishing indicators}
The following statement holds at arbitrary rank. For an object $V$, let
$I_{\rm low}(V)$ be the sum, with its multiplicities in $I(V)$, of those
simple constituents that occur in $I(1)$.

\begin{theorem}[Higher-indicator bound]\label{thm:higher-indicator}
Let $\mathcal C$ be a pseudo-unitary fusion category over $\CC$, with its
canonical spherical structure, and let $m\geq1$. If
$\nu_1(V)=\cdots=\nu_m(V)=0$, then
\[
 (m+1)\FPdim(I_{\rm low}(V))\leq D\FPdim(V).
\]
More strongly, the same inequality holds with $I_{\rm low}(V)$ replaced
by the full twist-one part of $I(V)$. Equality forces every other
constituent to have a nontrivial $(m+1)$st root of unity as twist.
In particular, if every simple constituent of $V$ is non-self-dual, then
$3\FPdim(I_{\rm low}(V))\leq D\FPdim(V)$.
\end{theorem}
\begin{proof}
Write $w_A=[I(V):A]\FPdim(A)\geq0$ and $T=D\FPdim(V)$.
The indicator identity \cite[Theorem 4.1]{NgSchauenburg} is
$\sum_Aw_A\theta_A^k=D\nu_k(V)$. Negative moments vanish by complex
conjugation. Expanding the nonnegative polynomial
\[
 |1+z+\cdots+z^m|^2=(m+1)+
 \sum_{k=1}^m(m+1-k)(z^k+z^{-k})\qquad(|z|=1)
\]
gives $\sum_Aw_A|1+\theta_A+\cdots+\theta_A^m|^2=(m+1)T$.
Every twist-one constituent contributes $(m+1)^2w_A$; discarding the
remaining nonnegative terms proves the bound. Equality requires the
polynomial to vanish at every remaining twist, proving its stated support.
The last assertion follows from additivity and
$\nu_1(X)=\nu_2(X)=0$ for a non-self-dual simple $X$.
\end{proof}
For $m=1$ this is the low-part estimate explicit in the proof of
\cite[Theorem 2.21]{OstrikRank3}. Its higher-indicator proof is included here. The arbitrary-sector
extension is recorded in Appendix~\ref{sec:twist-sectors}.

\subsection{The moment inequality}
For a nontrivial simple $X$, put $d=\FPdim(X)$ and $\epsilon=\nu_2(X)$.
The indicator formula \cite[Theorem 4.1]{NgSchauenburg} gives, with
$w_A=[I(X):A]\FPdim(A)\geq0$,
\begin{equation}\label{eq:three-moments}
 \sum_A w_A=Dd,\qquad \sum_A w_A\theta_A=0,\qquad
 \sum_A w_A\theta_A^2=D\epsilon.
\end{equation}
Here the twists are roots of unity, $\epsilon\in\{-1,1\}$ for a self-dual
simple and $\epsilon=0$ otherwise.

\begin{proposition}[Two-moment bound]\label{prop:optimized-moment}
If $L_X$ is any known contribution to the twist-one mass of $I(X)$, then
\[
 (3d+\epsilon)L_X\leq Dd(d+\epsilon).
\]
Consequently the following bounds require no choice of second indicator:
\begin{align}
 (3d+1)L_X&\leq Dd(d+1)&& (X\cong X^*),\label{eq:optimized}\\
 3L_X&\leq Dd&& (X\not\cong X^*).\label{eq:nonselfdual-third}
\end{align}
One may take $L_X=\FPdim(I_{\rm low}(X))=(LL^{\mathsf T}d)_X$.
\end{proposition}
\begin{proof}
For $P(z)=d+(d+\epsilon)z+dz^2$, expand $|P(z)|^2$ on the unit circle
and use \eqref{eq:three-moments}. This gives
\[
 \sum_Aw_A|P(\theta_A)|^2=Dd(3d+\epsilon)(d+\epsilon).
\]
The twist-one terms contribute at least $(3d+\epsilon)^2L_X$.
Since $3d+\epsilon>0$, division gives the assertion.
The function $e\mapsto(d+e)/(3d+e)$ is increasing on $[-1,1]$;
replace $\epsilon$ by one in the self-dual case.
\end{proof}

The familiar first-moment bound $2L_X\leq Dd$ and the weaker Fej\'er
bound $9L_X\leq D(3d+2)$ follow as well. The checker uses
\eqref{eq:optimized}--\eqref{eq:nonselfdual-third}; these suffice for
all the low-induction inequalities needed here. Positivity is essential:
these bounds are not applied to arbitrary spherical categories with signed
dimensions.

\subsection{A finite restriction test}
\begin{proposition}[Exhaustive low-induction certificate]\label{prop:low-certificate}
Fix the positive dimension character $d$. For each formal codegree $f$,
form the set of vectors $v\in\mathbb Z_{\geq0}^4$ satisfying
\begin{equation}\label{eq:column-search}
 v_0=1,\qquad \sum_i v_id_i=D/f,\qquad
 v_i\leq\min\{B_{0i},\lfloor\sqrt{B_{ii}}\rfloor\}\quad(i>0).
\end{equation}
Reserve $(1,0,0,0)^{\mathsf T}$ for the central unit. Enumerate all choices
of the other three columns, with the codegree multiplicities retained,
whose sum together with the unit column is $B_{\bullet0}$.
If every such choice fails the entrywise or positive-semidefinite condition
in \eqref{eq:induction-gram}, or one of
\eqref{eq:optimized}--\eqref{eq:nonselfdual-third}, then $R$ admits no
pseudo-unitary categorification.
\end{proposition}
\begin{proof}
A genuine restriction column satisfies its dimension equation, and
$M_{0A}=1$ on the four constituents of $I(1)$. The first row and the diagonal
of \eqref{eq:induction-gram} give the two bounds in
\eqref{eq:column-search}. Hence the actual low part occurs in the finite
list and satisfies every tested necessary condition.
\end{proof}

The computation is only linear algebra over a number field. Express the
four dimensions and $D/f$ in a rational power basis, row-reduce the dimension
equation, and enumerate its bounded free integer coordinates. When the
four dimensions are rationally independent, there is at most one possible
column for each $f$; a negative or nonintegral coefficient is already an
obstruction. When they are dependent, no uniqueness is presumed. Full
induction matrices, modular data and associators are never enumerated.

\subsection{The non-self-dual quadratic family}

For an integer $m\geq0$, let $S(m)$ be the commutative based ring with basis
$\{1,g,x,y\}$, involution $g^*=g$, $x^*=y$, and multiplication
\[
 g^2=1,\qquad gx=y,\qquad gy=x,\qquad
 x^2=y^2=g+m(x+y),\qquad xy=1+m(x+y).
\]

\begin{theorem}\label{thm:nonselfdual-quadratic}
For every $m\geq2$, the fusion ring $S(m)$ admits no fusion categorification
over $\mathbb C$.
\end{theorem}

\begin{proof}
Put $d=m+\sqrt{m^2+1}$. Thus $d^2=1+2md$, and the Frobenius--Perron
dimensions are $(1,1,d,d)$. By Schopieray's Galois-conjugacy theorem for
two-dimension fusion rings \cite[Lemma~3.1]{Schopieray}, any fusion
categorification would have a pseudo-unitary Galois conjugate with the same
fusion rules. It is consequently sufficient to exclude pseudo-unitary
categorifications.

Suppose one exists, and use its canonical spherical structure. Its dimension is
$D=2+2d^2=4+4md$. The four characters of $S(m)$, displayed by columns, are
\[
 \begin{pmatrix}
 1&1&1&1\\
 1&1&-1&-1\\
 d&-d^{-1}&i&-i\\
 d&-d^{-1}&-i&i
 \end{pmatrix}.
\]
Their formal codegrees are $D$, $2+2d^{-2}$, $4$, $4$, respectively.
The formal-codegree description of induction therefore gives
\[
 I(\mathbf1)\cong\mathbf1\oplus A\oplus B\oplus C,
 \qquad
 \operatorname{FPdim}(A)=d^2=1+2md,
 \quad
 \operatorname{FPdim}(B)=\operatorname{FPdim}(C)=1+md,
\]
where the four central simples are pairwise nonisomorphic and occur with
multiplicity one. Indeed, every irreducible representation of this commutative
rank-four ring is one-dimensional. Furthermore,
\[
 F I(\mathbf1)\cong
 \bigoplus_{T\in\operatorname{Irr}\mathcal C}T\otimes T^*
 \cong 4\mathbf1\oplus2mX\oplus2mY,
\]
where $X^*\cong Y$ correspond to $x,y$. By adjunction each of $F(A)$,
$F(B)$, $F(C)$ contains $\mathbf1$ once, and the displayed decomposition
shows that none contains $g$. Since $d$ is irrational, comparison of dimensions
now forces
\[
 [F(A):X]+[F(A):Y]=2m,\qquad
 [F(B):X]+[F(B):Y]=[F(C):X]+[F(C):Y]=m.
\]
These statements determine the sum of the two low-part dimensions without
requiring the individual restriction multiplicities:
\begin{align*}
 L&=\operatorname{FPdim}(I_{\mathrm{low}}(X))+
    \operatorname{FPdim}(I_{\mathrm{low}}(Y))\\
  &=2m(1+2md)+2m(1+md)=4m+6m^2d.
\end{align*}
Both $X$ and $Y$ are non-self-dual. Applying
\eqref{eq:nonselfdual-third} to each of them and adding gives $3L\leq2dD$.
But
\begin{align*}
 3L-2dD
 &=12m+18m^2d-2d(4+4md)\\
 &=12m+18m^2d-8d-8m(1+2md)\\
 &=4m+(2m^2-8)d>0
\end{align*}
for $m\geq2$, a contradiction.
\end{proof}


